\section{Experimental Protocol}
\label{sec:protocol}

\subsection{Mixed Evaluation}

The conventional baseline uses the frozen
\texttt{MIXED\_RANDOM\_V1} partition. Images from the available
source groups may occur across the development and test partitions,
while duplicate-control grouping is applied during split construction.
The resulting partition contains

\begin{itemize}
    \item 4,524 training images,
    \item 905 validation images, and
    \item 905 test images.
\end{itemize}

This protocol measures performance when training, validation, and test
samples are drawn from the same overall heterogeneous dataset.
Consequently, it represents the evaluation setting most likely to
benefit from source-specific regularities shared across partitions.

The validation partition is used for checkpoint selection and
training-control decisions. The test partition remains isolated until
the selected model is evaluated. CustomCNN\_V1 and MobileNetV2 are
first compared under this mixed protocol to establish their
within-dataset reference performance.

\subsection{Source-Disjoint Evaluation}

To evaluate sensitivity to image provenance, four frozen
source-disjoint composite folds are constructed. For each fold,
complete source groups are withheld from model development and used
only as the final test set. Images belonging to those held-out source
groups are absent from both training and validation.

Table~\ref{tab:source_disjoint_protocol} summarizes the four frozen
test configurations. Class counts are reported because source identity
and semantic class are partially entangled in the assembled dataset.

\begin{table*}[t]
\centering
\caption{Frozen Source-Disjoint Evaluation Protocols}
\label{tab:source_disjoint_protocol}
\small
\begin{tabular}{p{0.06\textwidth} p{0.43\textwidth} rrrr}
\hline
\textbf{Fold} &
\textbf{Held-Out Source Group(s)} &
\textbf{Desert} &
\textbf{Flooded} &
\textbf{Neither} &
\textbf{Total} \\
\hline

SD1 &
FDL UAV Flooded Areas + Aerial Landscape Images &
800 & 130 & 167 & 1,097 \\

SD2 &
Louisiana Flood 2016 + Satellite Image Classification &
1,131 & 102 & 161 & 1,394 \\

SD3 &
Disasters Dataset + Aerial Landscape Images &
800 & 1,493 & 137 & 2,430 \\

SD4 &
Cyclone/Wildfire/Flood/Earthquake +
Satellite Image Classification +
Forest Aerial/DeepGlobe-derived imagery &
1,131 & 613 & 1,000 & 2,744 \\

\hline
\end{tabular}
\end{table*}

The four folds intentionally represent different combinations of
source composition rather than four equivalent balanced test sets.
Accordingly, differences among SD1--SD4 are not interpreted as
controlled estimates of a single modality effect. Instead, they
measure robustness under different composite source-distribution
shifts.

CustomCNN\_V1 and MobileNetV2 are evaluated on all four folds using
the same preprocessing, class ordering, and evaluation metrics
defined in Section~\ref{sec:methodology}. The mean across SD1--SD4
is reported as a compact summary, while the individual folds remain
primary evidence because difficulty depends strongly on which source
groups are withheld.

Among the four baseline source-disjoint evaluations, SD1 produces
the lowest Macro-F1 for both principal architectures. SD1 is therefore
frozen as the focused robustness setting for the subsequent
domain-generalization experiments. The selection is based on the
baseline source-disjoint evaluation and is then kept fixed for all
later ERM, DG-AUG-V2, MixStyle, and RSC experiments.

\subsection{SD1 Development Partition}

For SD1, the complete FDL UAV Flooded Areas and Aerial Landscape
source groups form the 1,097-image test set. The remaining source
groups are separated into 4,488 training images and 749 validation
images:

\[
N_{\mathrm{train}}=4488,\qquad
N_{\mathrm{val}}=749,\qquad
N_{\mathrm{test}}=1097.
\]

The SD1 test set is not used for epoch selection, hyperparameter
selection, or checkpoint selection. Such decisions are based only on
the training and validation partitions.

The same frozen SD1 membership is reused across methods and random
seeds. This prevents differences in image membership from being
confounded with differences caused by initialization or training
strategy.

\subsection{Multi-Seed Robustness Evaluation}

Single-run comparisons can overstate the benefit of a robustness
method when optimization variability is substantial. The primary SD1
robustness comparison therefore evaluates three methods over three
random seeds:

\[
\mathcal{S}=\{42,123,2026\}.
\]

The repeated methods are

\[
\{
\mathrm{ERM},
\mathrm{DG\mbox{-}AUG\mbox{-}V2},
\mathrm{RSC}
\}.
\]

This design produces nine frozen models for the principal multi-seed
analysis. For every seed, data membership remains unchanged; only the
stochastic training trajectory changes. Results are summarized using
the arithmetic mean and sample standard deviation of Macro-F1.

MixStyle is not included in the three-seed comparison. Its late- and
mid-feature insertion variants are retained as diagnostic single-run
experiments and are interpreted separately rather than being presented
as seed-replicated DG evidence.

\subsection{External FloodNet UAV Evaluation}

The nine frozen multi-seed models are subsequently evaluated on the
independent labeled FloodNet UAV subset. The external set contains
398 images:

\begin{itemize}
    \item 51 Flooded scenes,
    \item 347 Non-Flooded scenes.
\end{itemize}

FloodNet Flooded maps to the internal Flooded class, whereas
Non-Flooded maps to Neither. No FloodNet image is assigned the Desert
ground-truth label.

The original three-class prediction vector is retained during external
inference. Desert predictions therefore remain incorrect outputs
instead of being removed before metric calculation. Supported-class
external Macro-F1 is computed over Flooded and Neither without
probability renormalization.

Before model evaluation, the external images are screened against all
6,334 primary-dataset images. Exact duplicate detection uses SHA-256
hashes. Perceptual-hash comparison is then used to identify potential
near-duplicate imagery that could survive resizing or re-encoding.
The frozen audit finds no exact overlap and no strongly suspicious
image-level match.

External results are reported independently for ERM, DG-AUG-V2, and
RSC across the same three seeds used in the SD1 robustness study.
This permits comparison of internal source-disjoint behavior with an
independently collected UAV domain.

\subsection{Consensus Failure Cohorts}

Predictions from all nine external models are used to define
agreement-based cohorts among the 347 true Neither FloodNet scenes.

Three groups are considered:

\begin{enumerate}
    \item \textbf{Consensus failure:} all nine models predict Flooded;
    \item \textbf{Consensus correct:} all nine models predict Neither;
    \item \textbf{Mixed agreement:} the nine predictions do not
    unanimously satisfy either condition.
\end{enumerate}

The frozen predictions yield 178 consensus failures,
56 consensus-correct Neither images, and 113 mixed-agreement images.
Cohort membership is determined before Grad-CAM or segmentation-mask
interpretation.

The unanimity criterion is intentionally strict. A consensus failure
must persist across variation in both training method and random seed,
making it more informative for studying recurring cross-model errors
than a mistake produced by only one checkpoint.

\subsection{Grad-CAM Sampling Protocol}

The Grad-CAM analysis uses three behavior categories:

\begin{enumerate}
    \item consensus-correct Flooded,
    \item consensus-correct Neither, and
    \item consensus Neither-to-Flooded failure.
\end{enumerate}

Within each category, examples are ranked using frozen prediction
confidence, and low-, median-, and high-confidence representatives are
selected algorithmically before heatmap interpretation. This reduces
the risk of manually selecting visually convenient examples after
seeing the explanations.

For every selected image, Grad-CAM is generated from the seed-42 ERM,
DG-AUG-V2, and RSC models. The resulting design contains

\[
3\ \text{behavior groups}
\times
3\ \text{confidence levels}
\times
3\ \text{methods}
=
27
\]

heatmaps.

The selected examples and resulting heatmaps are frozen before the
subsequent segmentation-based shortcut analysis.

\subsection{Post-Hoc Segmentation Audit}

The segmentation audit focuses on the two unanimous Neither cohorts:

\begin{itemize}
    \item 178 consensus Neither-to-Flooded failures,
    \item 56 consensus-correct Neither scenes.
\end{itemize}

Official FloodNet semantic masks provide the scene-composition labels.
For each of the ten segmentation classes, two statistical comparisons
are performed.

First, semantic-class \emph{presence} is represented as a binary
image-level variable and compared between cohorts using Fisher's exact
test. Second, semantic-class \emph{pixel coverage} is represented by
the fraction of pixels assigned to the class and compared using the
Mann--Whitney \(U\) test.

Because ten semantic categories are examined, family-wise
multiple-testing correction is performed using the Holm procedure.
Both raw and Holm-adjusted \(p\)-values are retained in the evidence
tables, while adjusted values are used for the principal significance
interpretation.

The segmentation audit is post hoc and exploratory. It is intended to
identify scene characteristics statistically associated with systematic
external failure, not to establish a causal mechanism for model
predictions.

\subsection{Experiment-Freezing and Model-Selection Rules}

The experimental design follows several rules intended to reduce
leakage and retrospective selection.

First, model checkpoints are selected using validation loss rather than
test performance. Second, membership of the mixed and source-disjoint
partitions remains fixed during architecture and method comparisons.
Third, multi-seed conclusions are based on the complete prespecified
seed set rather than a favorable individual run.

Primary result tables are generated from frozen experiment outputs.
Diagnostic experiments that were superseded or were not replicated
across the required seed set remain identifiable as diagnostic
evidence and are not promoted to the same evidentiary status as the
primary comparisons.
